
\section{The deformation variety for the figure eight knot complement}
\label{computations}

We fix the usual ideal $2$-tetrahedra triangulation of the figure eight knot
complement (which we refer as the standard triangulation). In this section we
prove the following theorem.

\begin{theorem}\label{3components}
  Given the standard triangulation of the figure eight knot complement, the
  deformation variety $\tmop{Defor} ( M, \mathcal{T} )$ is the union of 3
  distinct smooth affine (irreducible algebraic) varieties of dimension 2 and
  is connected.
\end{theorem}

As said in the introduction, we strongly believe that the character variety of the figure 
eight knot complement does not contain any other irreducible component with irreducible 
representations. A proof should appear elsewhere.

In order to prove the theorem, we will compute 3 affine varieties $D_{1}
,D_{2} ,D_{3}$, 3 birational maps $\pi_{1} , \pi_{2} , \pi_{3}$ defined
everywhere on the $D_{i}$'s and $3$ polynomials $H_{1} ,H_{2} ,H_{3}$ in
$\mathbb{Q} [ X_{1} ,X_{2} ,X_{3} ]$ such that :
\begin{itemize}
  \item (a) $\tmop{Defor} ( M, \mathcal{T} ) = \cup_{i=1 \ldots 3} D_{i}$.
  
  \item (b) $V ( H_{i} ) = \overline{\pi_{i} ( D_{i} )} ,i=1 \ldots 3$.
  
  \item (c) $\pi_{i} ( D_{i} ) \subset V ( H_{i} )$ is smooth and $\pi_1$ realizes a homemorphism on its image.  
  \item (d) $V ( H_{i} )$ is irreducible in $\mathbb{C}^{3}$  or equivalently $H_{i}$ is
  irreducible in $\mathbb{C} [ X_{1} ,X_{2} ,X_{3} ]$.
  
  \item (e) $D_{i} \neq D_{j}$ \ for $i \neq j$ and $i,j=1 \ldots 3$.
    \item (f) $\cap_{i=1 \ldots 3} D_{i} \neq \emptyset$.
\end{itemize}
Under assumptions (a) to (f), $\pi_{i} ( D_{i} )$ is a  smooth Zariski-open subset
in the affine irreducible algebraic variety $V ( H_{i} )$, it is thus
connected -- see \cite[4.16]{Mumford}. As $\pi_{i}$ is  is a homeomorphism on its image, $D_{i}$ is also smooth and
connected. As $\cap_{i=1 \ldots 3} D_{i} \neq \emptyset$, then $\tmop{Defor}(M,\mathcal T)$
is also connected. So the theorem follows from those properties.

Let us describe how to compute $D_{i} , \pi_{i} ,H_{i} ,i=1 \ldots 3$.

\subsection{Computational tools}

The main computational tool for computing these objects is Gr\"obner basis. We recall briefly 
that a Gr\"obner basis of an ideal $\mathcal I$ in $\mathbb Q[X_1,\ldots,X_n]$ is a set of 
generators of $\mathcal I$ such that, for any polynomial $P$, there is a unique preferred 
polynomial, denoted by $P\,\; \mathrm{mod}\, \;\mathcal I$, congruent to $P$ modulo $\mathcal I$. 
This polynomial is called the \emph{normalform}. For more details, we refer to \cite{CLO1997}. 
Note that such a basis is uniquely defined once chosen an admissible ordering on the monomials. Once a Gröbner basis is known, one can compute the associated Hilbert polynomial and then the (Hilbert) dimension and the (Hilbert) degree \cite[Section 9.3]{CLO1997}.

An interesting ordering is what is called \emph{elimination ordering}. If $\mathcal I$ is seen 
as an ideal of $\mathbb Q[Y_1,\ldots, Y_k][X_1,\ldots,X_r]$,  an elimination ordering verifies 
$Y_i<X_j$ for all $i$, $j$. It allows a straightforward computation of a Gröbner basis for $$
\mathcal I'=\mathcal I\cap \mathbb Q[Y_1,\ldots,Y_k],$$ i.e. elimination of variables \cite[Ch. 
3]{CLO1997}. Indeed, the set of zeroes of $V(\mathcal I')$ is the Zariski-closure of the 
projection of $V(\mathcal I)$ on the first variables \cite[Section 4.4]{CLO1997}.

We will often eliminate some variables, especially when we may express a variable by a 
rational expression in the other ones. But, the projections might create spurious components when 
taking the Zariski closures: we work with constructible sets rather than algebraic varieties. As our computations are on the edge of what can be actually done, 
we want to avoid these components. So we will make a frequent use of saturation. Indeed, for 
an ideal $\mathcal I$ in $\mathbb Q[X_1,\ldots,X_n]$ and a  polynomial
$f$, the saturation of $\mathcal I$ by $f$ is the ideal $\mathcal I:f^\infty$,
which consist in the set of polynomials $P$ such that for some $k$,
the polynomial $Pf^k$ belongs to $\mathcal I$. Geometrically, when
$\mathcal I$, the zeroes set of $\mathcal I:f^\infty$ is the Zariski closure of $V(\mathcal I)\setminus V(f)$. The ideal $\mathcal I:f^\infty$ is computed by taking a Gröbner basis of $(\mathcal I+<tf-1>)\cap \mathbb Q[X_1,\ldots ,X_n]$, where $t$ is an independant variable -- thanks to an elimination ordering.
 
Let us formalize the operations we need and explain their scope. Here is a list of algorithms we will use later on:
\begin{itemize}
  \item [$\alpha$] Computing an ideal, in the sense of computing a Gröbner basis.
  
  \item [$\beta$] Saturating an ideal $I \in \mathbb{Q} [ X_{1} , \ldots
  ,X_{n} ]$ by  $f \in \mathbb{Q} [ X_{1} , \ldots ,X_{n}
  ]$, in the sense of computing a Gröbner basis of the ideal $I:f^{\infty}$. We
  naturally extend this algorithm for saturating an ideal $I$ by a set of
  polynomials $\{ f_{1} , \ldots ,f_{l} \}$ in order to compute $I:
  \left( \prod_{i=1 \ldots l} f_{i} \right)^{\infty}$. This is simply achieved by iteratively 
  saturating by $f_1$, then $f_2$, etc \cite[Ch. 4, prop. 10]{CLO1997}.
  
  \item [$\gamma$] Factorizing a polynomial $p \in \mathbb{Q} [ X_{1} ,
  \ldots ,X_{n} ]$ in $\mathbb{Q} [ X_{1} , \ldots ,X_{n} ]$.
  

  \item [$\eta$] \tmtextbf{for small systems (low degree, few number of
  variables)} Computing the prime decomposition of an ideal generated by the
  equations, removing inclusions. For those who are not familar with such
  objects, just retain that this means computing the decomposition of an
  algebraic variety in $\mathbb{Q}$- irreducible  (defined by prime
  ideals in $\mathbb{Q} [ X_{1} , \ldots ,X_{n} ]$) and non redundant
  components.
  
  \item [$\theta$] \tmtextbf{ Not usually implemented, see later on.} Testing if a polynomial $p \in \mathbb{Q} [ X_{1} ,X_{2}
  ,X_{3} ]$ is irreducible in $\mathbb{C} [ X_{1} ,X_{2} ,X_{3} ]$ or not.
\end{itemize}
 
Note that the four first algorithms are implemented in the usual computer algebra systems, such 
as Maple or Sage. Sadly, our problem seems far beyond their scope. So a great deal of our 
method is to use the shape of our equations to simplify the problem and getting to the point we can use these routines.

However the last algorithm [$\theta$] is special: it is not implemented in usual systems -- which 
work with number fields. Moreover, it is only used at the very last step of our proof of the 
theorem. So we will be more precise later on.
 

\subsection{An algebraic representation of $\tmop{Defor} ( M, \mathcal{T} )$}\label{ss:equations}

Given $M$ and $\mathcal{T}$ a triangulation of $M$ with $n$ tetrahedra, the
deformation variety $\tmop{Defor} (M, \mathcal{T} )$ is defined in
$\mathbb{C}^{16n}$ by a system of internal relations for each tetrahedron and
gluing equations for adjacent tetrahedra. Moreover, using the internal
relations, we may directly express the face parameters in terms of the edge
parameters. From now on, we consider $\tmop{Defor} (M, \mathcal{T} )$ as a
subset of $\mathbb{C}^{12n}$.

The system, more precisely, is given in terms of variables $Z= (z_{ij} (T))
\in \mathbb{C}^{12n}$, where $T$ denotes a tetrahedron in the triangulation
(containing $n$ tetrahedra) and $ij$ one of its oriented edges.
\begin{itemize}
  \item $L_{e}$ (resp. $L_{f}$) denotes the set of the $2n$ edge (resp. face)
  equations;
  
  \item $L_{c}$ denotes the polynomials defining the internal relations (also
  called cross-ratio relations), of the form $z_{ik} (T)  (1-z_{ij} (T)) -1$
  and $z_{ij} (T)  (1-z_{il} (T)) -1$ for each tetrahedron;
\end{itemize}
In this terms, $\tmop{Defor} (M, \mathcal{T} )$ is the algebraic variety:
\[ \tmop{Defor} (M, \mathcal{T} ) =V ( \langle L_{e} \cup L_{f} \cup L_{c}
   \rangle ) = \{{Z} \in \mathbb{C}^{12n} |P(Z)=0,P \in L_{e} \cup
   L_{f} \cup L_{c} \} . \]
Note that $\tmop{Defor} (M, \mathcal{T} ) \subset (\mathbb{C} \setminus
\{0,1\})^{12n}$, because of the internal relations $L_{c}$. This remark will be important later on.

We now display the equations for the figure eight knot complement. The triangulated structure (see section \ref{section:intro-eight}) associated
with the figure eight knot complement is made of two tetrahedra decorated by
the coordinates $(u_{ij} )_{i \neq j \in \{1, \ldots ,4\}}$ (resp. $(v_{ij}
)_{i \neq j \in \{1, \ldots ,4\}}$) with cross-ratio relations defined as the
roots of the system $L_{c_{1}} =0$ (resp. $L_{c_{2}} =0$) :
\begin{center}
  $L_{c_{1}} = \left\{ \begin{array}{c}
    u_{12} u_{13} -u_{13} +1\\
    u_{12} u_{14} -u_{12} +1\\
    u_{21} u_{23} -u_{21} +1\\
    u_{21} u_{24} -u_{24} +1\\
    u_{31} u_{34} -u_{31} +1\\
    u_{32} u_{34} -u_{34} +1\\
    u_{41} u_{43} -u_{43} +1\\
    u_{42} u_{43} -u_{42} +1
  \end{array} \right. \begin{array}{c}
    \\
    \\
    
  \end{array}$ , \ $L_{c_{2}} = \left\{ \begin{array}{c}
    v_{12} v_{13} -v_{13} +1\\
    v_{12} v_{14} -v_{12} +1\\
    v_{21} v_{23} -v_{21} +1\\
    v_{21} v_{24} -v_{24} +1\\
    v_{31} v_{34} -v_{31} +1\\
    v_{32} v_{34} -v_{34} +1\\
    v_{41} v_{43} -v_{43} +1\\
    v_{42} v_{43} -v_{42} +1
  \end{array} \right.$ ,
\end{center}
and gluing edge (resp. face) relations defined as the roots of the system
$L_{e} =0$ (resp. $L_{f} =0$) :
\begin{center}
  $L_{e} = \left\{ \begin{array}{c}
    u_{13} u_{14} u_{23} v_{21} v_{24} v_{31} -1\\
    u_{31} u_{32} u_{41} v_{12} v_{13} v_{42} -1\\
    u_{21} u_{24} u_{34} v_{23} v_{41} v_{43} -1\\
    u_{12} u_{42} u_{43} v_{14} v_{32} v_{34} -1
  \end{array} \right.$ , $L_{f} = \left\{ \begin{array}{c}
    u_{21} u_{31} u_{41} v_{21} v_{31} v_{41} -1_{\nosymbol}\\
    u_{12} u_{32} u_{42} v_{13} v_{23} v_{43} -1\\
    u_{13} u_{23} u_{43} v_{12} v_{32} v_{42} -1\\
    u_{14} u_{24} u_{34} v_{14} v_{24} v_{34} -1
  \end{array} \right.$ .
\end{center}
We set $Z= \{ u_{ij} ,v_{i,j} ,i \neq j \in \{ 1, \ldots ,4 \} \}$.


\subsection{Overview of the computations.}

Theoretically, the theorem could be proven using a prime decomposition of the ideal 
$$I_D=\langle L_{e} \cup L_{f}\cup L_{c} \rangle.$$
Indeed, as $\mathrm{Defor}(M, \mathcal T)$ is the zeroes set of this ideal, a prime decomposition would give the three components and we would be able to check their properties.
However, our system depends on $12n$ variables (24 variables in the
case of the figure eight knot complement) and is of rather high degree
(degree $6$ for some equations for the figure eight knot complement). So we are far beyond the scope of the state of the art algorithms of prime decomposition.

We will use the shape of the equations to simplify the system. The method we present is rather general for this kind of systems up to some points (number of factors
in some polynomials). It could be applied, in principle on numerous systems
with the same properties. For example, it can replace the one used in \cite{FKR14}.
Before explaining our method, let us define one more object: \emph{forbidden points}. 

Denote by $S_{D}=L_{e} \cup L_{f}\cup L_{c}$ the set of polynomials generating the ideal $I_{D}$.
Note that if $z \in V ( I_{D} )\subset \bC^{12n}$, then none of its coordinates
  belong to $\{ 1,0 \}$ \ because of the equations induced by the
  {\emph{internal relations}} $L_c$. We hence call {\tmem{forbidden
  points}} the points with any coordinate in $\{ 0,1 \}$ and we introduce 
  $F_{D}= \{ z,z-1,z \in Z \} \subset \mathbb{Q} [ Z ]$. It may seem useless at this point as this is contained in the internal relations.
  But recall that we will eliminate a lot of variables through projections. During these projections, as said before, some spurious components (i.e. all included in forbidden points) may appear. So we will always keep the equations defining the forbidden points and saturate our ideals by these equations. This will lower the number of components and ease the computations.
  
 Here are the steps of our method: 
 
  \begin{itemize}
\item {\bf Step A: elimination.} 
We iteratively detect equations in $S_{D}$ which are affine in a variable, i.e.  of the form 
$$D_{z} ( X )  z_{} -N_{z} ( X )\textrm{ with }X \subset Z \setminus \{ z \},$$ 
and with the additional condition that $D_{z} ( X )$
  cancels only at forbidden points. We can then substitute $z$ by $\frac{N_{z} ( X )}{D_{z}
  ( X )}$ in the other polynomials generating $I_{D}$, take the numerators of
  these equations and remove their factors (algorithm [$\gamma$]) that cancel
  exclusively at forbidden points (and thus never cancel on $\tmop{Defor} ( M,
  \mathcal{T} )$). This gives a new ideal with one less variable and one less equation.
  
At the end of the process, denote by $S_{Y} \subset \mathbb{Q} [ Y ]$ with $Y \subset Z$ the
  remaining equations. We keep in mind, through the set $R_{Y}
  \subset \mathbb{Q} [ Z ]$, the equations used for the substitutions. Of course, by substituting in $R_Y$, we may assume that all these equations express a substituted variable in terms of the variables in $Y$, i.e. they have the form: $D_{z} ( Y ) z-N_{z} ( Y )$, with $z
  \in Z \setminus Y$.
  
We also want to keep in mind the forbidden points. So denote by $F_{Y} \in \mathbb{Q} [ Y ]$
 the list of prime factors that appear in $F_{D}$ after having performed the same substitution.
  
We can interpret the relations $R_{Y} =0$ as a projection $\pi_{Y} : \tmop{Defor} (
  M, \mathcal{T} ) \rightarrow V ( \langle S_{Y} \rangle )$ which is not
  necessarily surjective. However, the complement of the image $V ( \langle S_{Y}
  \rangle ) \setminus \pi_{Y} ( \tmop{Defor} ( M, \mathcal{T} ) )$ only contains
  forbidden points which might have been introduced by the use of
  rational fractions in the substitutions.
  
  So we compute (algorithm $[ \beta ]$) the ideal $I_{Y}$, which is the saturation of $\langle S_{Y}  \rangle$ by $F_{Y}$ in order to get $V ( I_{Y} ) = \overline{\pi_{Y} (
  \tmop{Defor} )}$. The union of the zeroes of the polynomials $F_{Y}$ defines
  the {\tmem{forbidden points}} of $V ( I_{Y} )$ which are images by $\pi_{Y}$
  of the forbidden points of $\mathbb{C}^{12n}$.
  
Putting this in a diagram, we have :
\begin{center}
  \begin{tabular}{ccc}
    $\mathbb{C}^{12n}$ & $\rightarrow$ & $\mathbb{C}^{\sharp Y}$\\
    $\tmop{Defor} ( M, \mathcal{T} )$ & $\approx_{\pi_{Y}}$ & $V ( I_{Y} )
    \setminus ( \cup_{f \in F_{Y}} V ( f ) )$\\
    $( x,y )$ & $\rightarrow$ & $y$\\
    $\left( \frac{N_{x} ( y )}{D_{x} ( y )} ,y \right)$ & $\leftarrow$ & $y.$
  \end{tabular}
\end{center}
  \item {\bf Step B: splitting.} In $I_{Y}$, we find a polynomial
  $P_{Y} \in \mathbb{Q} [ Y ]$ that factorizes (algorithm [$\gamma$]) into
  $k$ factors in $\mathbb{Q} [ Y ]$ : $P_{Y} = \prod_{i=1 \ldots k} P_{Y,i}$.
  
  By construction, as $\tmop{Defor} ( M, \mathcal{T} )$ does not contain
  forbidden points, we then get $\tmop{Defor} ( M, \mathcal{T} ) = \cup_{i=1
  \ldots k} V ( J_{Y,i} )$, where $J_{Y,i} \subset \mathbb{Q} [ Z ]$ is the
  saturation of $I_{D} + \langle P_{Y,i} \rangle$ by $F_{D}$.
  
  We can even further saturate our ideals, in order to avoid the same component appearing twice. We  perform cascading
  saturations  for some permutation $\sigma$ of the indices: we iteratively compute
  (algorithm $[ \beta ]$)  the ideal $I_{Y, \sigma ( i )} =I_{Y} + \langle
  P_{Y, \sigma ( i )} \rangle$ and saturate this ideal by $F_{Y}$ as well as
  with $\{ P_{Y, \sigma ( 1 )} , \ldots ,P_{Y, \sigma ( i-1 )} \}$ \ in order
  to remove possible components that are already in $I_{Y, \sigma ( 1 )} ,
  \ldots ,I_{Y, \sigma ( i-1 )}$. The efficiency of the process strongly
  depends on the chosen permutation.
  
At the end of this step, one can set $D_{i} =V ( I_{Y,i} +I_{D} )$. By construction, $D_i$ is the inverse image by $\pi_i$ of the zeroes set of $I_{Y,i}$. And we get the decomposition: 
$$\tmop{Defor} ( M,
  \mathcal{T} ) = \cup_{i=1 \ldots k} D_{i}.$$
  
  \item {\bf Step C: further elimination.} The ideals $I_{Y,i}$ are still too big to perform 
  directly a prime decomposition. So, for each $I_{Y,i}$, we apply once again step A. We may 
  even try to iteratively apply Steps A and B and so on, but in our case, a single additional 
  pass of Step A turns out to be enough. So we explain only this single pass.
  
  Indeed, for each $I_{Y,i}$, we find a polynomial that is affine in
  one of the variables and, moreover, can be written $a_{i} ( z ) y+b_{i} ( z
  )$ with $z \in Y_{i} =Y \setminus \{ y \}$ and $a_{i} ( z^{} )$ only vanishing at forbidden points. We stress that 
  the suitable variable $y$ to eliminate depends on the different ideals $I_{Y,i}$.
  
  The result of Step A is then an ideal $I_{Y_{i}} \subset \mathbb{Q}
  [ Y_{i} ]$ together with a projection $\pi_{Y_{i}} :V ( I_{Y,i} ) \rightarrow V (
  I_{Y_{i}} )$ and a set of polynomials
  $F_{Y_{i}}$ (defining the {\tmem{forbidden points}} of $V ( I_{Y_{i}} )$)
  such that we get: $$\overline{\pi_{Y_{i}} ( V ( I_{Y,i} ) )} =V ( I_{Y_{i}} ).$$
  
  Setting $\pi_{i} = \pi_{Y_{i}} \circ \pi_{Y}$, we may write
  $\overline{\pi_{i} ( D_{i} )} =V ( I_{Y_{i}} )$ and get the diagram similar to the previous 
  one:
 
\begin{center}
    \begin{tabular}{ccc}
      $\mathbb{C}^{12n}$ & $\rightarrow$ & $\mathbb{C}^{\sharp Y-1}$ \\

      $D_{i}$ & $\approx_{\pi_{Y_i}}$ & $V ( I_{Y_i} ) \setminus ( \cup_{f \in
      F_{Y_i}} V ( f ) )$ \\
%
%      $( x,y,z )$ & $\rightarrow$ & $( y,z )$\\
%
%      $\left( \frac{N_{x} \left( \frac{N_{y,i} ( z )}{D_{y,i} ( z )}
%      \right)}{D_{x} \left( \frac{N_{y,i} ( z )}{D_{y,i} ( z )} \right)} ,
%      \frac{N_{y,i} ( z )}{D_{y,i} ( z )} ,z \right)$ & $\leftarrow$ & $\left(
%      \frac{N_{y,i} ( z )}{D_{y,i} ( z )} ,z \right)$
    \end{tabular}
\end{center}
%and 
%\begin{center}
%    \begin{tabular}{ccc}
%     $\mathbb{C}^{\sharp Y}$ &
%      $\rightarrow$ & $\mathbb{C}^{\sharp Y-1}$\\
%
%      $V ( I_{Y,i} ) \setminus ( \cup_{f \in
%      F_{Y}} V ( f ) )$ & $\approx_{\pi_{Y_{i} ,i}}$ & $V ( I_{Y_{i}} )
%      \setminus \left( \cup_{f \in F_{Y_{i}}} V ( f ) \right)$\\
%
%      $( y,z )$ & $\rightarrow$ & $z$\\
%
%      $\left(
%      \frac{b_{i} ( z )}{a_{i} ( z )} ,z \right)$ & $\leftarrow$ & $z$
%    \end{tabular}
%\end{center}
%  with $\tmop{Defor} ( M, \mathcal{T} ) = \cup_{i=1 \ldots k} D_{i}$,
%  $\overline{\pi_{Y} ( D_{i} )} =V ( I_{Y,i} )$ and $\overline{\pi_{Y_{i}
%  ,i} ( V ( I_{Y,i} ) )} =V ( I_{Y_{i} ,i} )$, $\overline{\pi_{i} ( D_{i} )} =
%  \overline{\pi_{Y_{i} ,i} \circ \pi_{Y} ( D_{i} )} =V ( I_{Y_{i} ,i} )$.
Beware that the different projections $\pi_{Y_i}$ are only defined on their respective component $D_i$ and do not land in the same space, as the last eliminated variable may differ for different components.
  
  \item {\bf Step D: Prime decomposition.} At this stage, each ideal $I_{Y_{i}}$ belongs to 
  the scope of a Prime Decomposition algorithm (named here algorithm [$\eta$]). Doing so, we 
  extract the $\mathbb{Q}$-irreducible components or just a check that the $I_{Y_{i}}$ 
  is prime. We denote by $I_{Y_i,k}$ the different prime ideals appearing in the decomposition of $I_{Y_i}$. 
\end{itemize}


The Steps $A$ to $D$ are general functions that also have been applied
for other computations, for example for computing unipotent solutions for many
cases of varieties with a triangulation involving up to $4$ tetrahedra (\cite{FKR14}).

At this point, we have a decomposition of $\mathrm{Defor}(M,\mathcal T)$ in components, for which we have a prime decomposition. Proving the theorem for this decomposition is now just a matter of checking properties and here the method is very specific to the figure eight knot complement. Other manifolds $M$ may give other behaviors. In particular, it turns out in our case that, at this step, each ideal $I_{Y_i}$ has a component principal and prime and maybe another component of lesser dimension and which is redundant:
\begin{itemize}
  \item {\bf Specific Step $1$: irreducibility and smoothness} Check that for each $i$ there is a unique prime ideal of maximal dimension among the prime factors. We denote here by $I_{Y_i,1}$ this prime factor. Check that $\mathrm{Defor}(M,\mathcal T)$ is included in the union of the $\pi_i^{-1}(I_{Y_i,1})$. In other terms, the components of lesser dimension are redundant. Then check that each $I_{Y_{i} ,1} \subset \mathbb{Q} [ Y_{i} ]$ is  a  principal ideal of $\mathbb{Q} [ Y_{i} ]$ (the output of algorithm
  [$\alpha$] is a unique polynomial), and, if so, denote by $H_{i} \in
  \mathbb{Q} [ Y_{i} ]$ its generator.
  
  Eventually verify (algorithm [$\gamma$]) that $H_{i}$ is irreducible in
  $\mathbb{C} [ Y_{i} ]$ and check that the singular points of each $H_{i}$
  are exclusively forbidden points.
  
  \item {\bf Specific Step $2$: connectedness.} Let $I_{i}$ be the ideal generated by $I_{Y_i,1}$ and $I_D$. It verifies that $V ( I_{i} ) =D_{i}$. We compute (algorithm [$\alpha$]) the ideal $I_{1} + \ldots +I_{k}$. And we check if is not trivial.
\end{itemize}
If the results of these two specific steps are indeed what is announced, then the theorem is easily proven: $\pi_{i} ( D_{i} )$ is a dense open
  set in $V ( H_{i} )$ which is a smooth affine (irreducible) variety so that
  $\pi_{i} ( D_{i} )$ and thus $D_{i}$ are smooth and connected. Moreover,  as the ideal $I_{1} + \ldots +I_{k}$ is not trivial, we get that $D_{1} \cap \ldots \cap D_{k}$ is a non empty 
  set of points. As each $D_i$ is connected, we conclude that $\tmop{Defor} ( M, \mathcal{T} ) = \cup_{i=1 \ldots k} D_{i}$ is connected. 



\subsection{Explicit computations for the Figure Eight Knot complement.}\label{results}

We now review the different steps in the case of the figure eight knot complement. We indicate the choices that can be done to ensure the completion of our method. Beware that the feasibility of the computations depends on these choices. For example, there are choices of the eliminated variables in Step A that lead to unfeasible computations. Let us mention that the reader may find the Maple worksheet where this computation is implemented: see \cite{Worksheet}.

\subsubsection{{\bf Step A}}

We apply step A on $S_{D} =L_{c_{1}} \cup L_{c_{2}} \cup L_{e} \cup
L_{f}$. With the set $Y= \{ v_{1,4} ,u_{4,3} ,v_{4,3} ,v_{3,4} \}$, one get the following expressions for $\pi_Y$:
{\small \begin{center}
  $\pi_{Y} : \left\{ \begin{array}{l}
    u_{12} =1- \frac{1}{u_{13}} ,u_{13} =1- \frac{1}{u_{14}} ,u_{14} =1-
    \frac{(1-v_{34} )(1-u_{24} )}{1-v_{24}} ,u_{21} = \frac{1}{1-u_{23}} ,\\
    u_{23} = \frac{1}{1-u_{24}} ,u_{24} =1+v_{14} (1-v_{24} )
    \frac{u_{43}}{1-u_{43}} ,u_{31} = \frac{1}{1-u_{34}} ,u_{32} =1-
    \frac{1}{u_{34}} ,\\
    u_{34} =- \frac{1-u_{43}}{u_{43} v_{14} (1-v_{43} )} ,u_{41} =1-
    \frac{1}{u_{43}} ,u_{42} = \frac{1}{1-u_{43}} ,\\
    v_{12} = \frac{1}{1-v_{14}} ,v_{13} =1- \frac{1}{v_{14}} ,v_{21} =1-
    \frac{1}{v_{24}} ,v_{23} = \frac{1}{1-v_{24}} ,\\
    v_{24} =1- \frac{v_{43} (1-v_{14} )}{u_{43} v_{14}^{2} (1-v_{34} )} -
    \frac{v_{43} (1-v_{43} )(1-v_{14} )}{(1-u_{43} )v_{14} (1-v_{34} )} ,\\
    v_{31} = \frac{1}{1-v_{34}} ,v_{32} =1- \frac{1}{v_{34}} ,v_{41} =1-
    \frac{1}{v_{43}} ,v_{42} = \frac{1}{1-v_{43}} .
  \end{array} \right.$
\end{center}}
Moreover the ideal $I_{Y}$ is generated by $3$ polynomials in $\mathbb{Q} [
Y ]$ of maximal degree $5$ in each variable, which are too large to be printed
here. Neither do we print $F_{Y}$ which is quite large and easy to compute using a
simple substitution and some factorizations.

Note that $\pi_{Y}$ is not simplified here, the expression having only
variables of $Y$ in the left sizes of the equations being a little too large
to be printed in this article.

\subsubsection{{\bf Step B}}

We are looking for a polynomial in $I_Y$ that may be factorized. As said before, $I_Y$ is generated by $3$ polynomials. Computing the gcd in $\mathbb{Q} [ v_{1,4} ,u_{4,3} ,v_{4,3} ]$ of all the resultants of 2 of these $3$ polynomials wrt $v_{3,4}$ we then obtain a
polynomial that belongs to $I_{Y}$ and factorizes into $3$ factors. We use this polynomial as $P_{Y}$ and get the three factors:
\begin{itemize}
  \item {\small{$P_{Y,1} =-u_{4,3} v_{1,4}^{2} v_{4,3}^{2} +u_{4,3}^{2} v_{1,4}
  v_{4,3} +u_{4,3} v_{1,4}^{2} v_{4,3} +u_{4,3} v_{1,4} v_{4,3}^{2}
  -u_{4,3}^{2} v_{1,4} -3 \hspace{0.17em} u_{4,3} v_{1,4} v_{4,3} +v_{1,4}
  u_{4,3} +u_{4,3} v_{4,3} +v_{1,4} v_{4,3} -v_{4,3}$}}
  
  \item {\small{$P_{Y,2} =-v_{1,4}^{2} v_{4,3}^{2} +2 \hspace{0.17em}
  v_{1,4}^{2} v_{4,3} +v_{1,4} v_{4,3}^{2} +u_{4,3} v_{4,3} -v_{1,4}^{2}
  -v_{1,4} v_{4,3} -v_{4,3}$}}
  
  \item {\small{$P_{Y,3} =u_{4,3}^{5} v_{1,4}^{4} v_{4,3}^{2} +u_{4,3}^{4}
  v_{1,4}^{4} v_{4,3}^{3} +u_{4,3}^{3} v_{1,4}^{4} v_{4,3}^{4} -u_{4,3}^{5}
  v_{1,4}^{4} v_{4,3} -3 \hspace{0.17em} u_{4,3}^{4} v_{1,4}^{4} v_{4,3}^{2}
  -u_{4,3}^{4} v_{1,4}^{3} v_{4,3}^{3} -3 \hspace{0.17em} u_{4,3}^{3}
  v_{1,4}^{4} v_{4,3}^{3} -2 \hspace{0.17em} v_{1,4}^{3} u_{4,3}^{3}
  v_{4,3}^{4} +2 \hspace{0.17em} v_{1,4}^{4} u_{4,3}^{4} v_{4,3} +2
  \hspace{0.17em} v_{1,4}^{3} u_{4,3}^{4} v_{4,3}^{2} +3 \hspace{0.17em}
  u_{4,3}^{3} v_{4,3}^{2} v_{1,4}^{4} +6 \hspace{0.17em} v_{1,4}^{3}
  u_{4,3}^{3} v_{4,3}^{3} +u_{4,3}^{3} v_{1,4}^{2} v_{4,3}^{4} -u_{4,3}^{4}
  v_{1,4}^{3} v_{4,3} -3 \hspace{0.17em} u_{4,3}^{4} v_{1,4}^{2} v_{4,3}^{2}
  -u_{4,3}^{3} v_{1,4}^{4} v_{4,3} -6 \hspace{0.17em} u_{4,3}^{3} v_{1,4}^{3}
  v_{4,3}^{2} -5 \hspace{0.17em} u_{4,3}^{3} v_{1,4}^{2} v_{4,3}^{3}
  -u_{4,3}^{2} v_{1,4}^{3} v_{4,3}^{3} +3 \hspace{0.17em} v_{1,4}^{2}
  u_{4,3}^{4} v_{4,3} +2 \hspace{0.17em} v_{1,4}^{3} u_{4,3}^{3} v_{4,3} +10
  \hspace{0.17em} u_{4,3}^{3} v_{4,3}^{2} v_{1,4}^{2} +2 \hspace{0.17em}
  u_{4,3}^{3} v_{1,4} v_{4,3}^{3} +2 \hspace{0.17em} u_{4,3}^{2} v_{1,4}^{3}
  v_{4,3}^{2} +3 \hspace{0.17em} u_{4,3}^{2} v_{1,4}^{2} v_{4,3}^{3}
  -u_{4,3}^{4} v_{1,4} v_{4,3} -6 \hspace{0.17em} v_{1,4}^{2} u_{4,3}^{3}
  v_{4,3} -2 \hspace{0.17em} u_{4,3}^{3} v_{1,4} v_{4,3}^{2} -v_{1,4}^{3}
  u_{4,3}^{2} v_{4,3} -6 \hspace{0.17em} v_{1,4}^{2} u_{4,3}^{2} v_{4,3}^{2}
  -2 \hspace{0.17em} u_{4,3}^{2} v_{1,4} v_{4,3}^{3} +2 \hspace{0.17em}
  u_{4,3}^{3} v_{1,4} v_{4,3} +u_{4,3}^{3} v_{4,3}^{2} +3 \hspace{0.17em}
  u_{4,3}^{2} v_{1,4}^{2} v_{4,3} +2 \hspace{0.17em} u_{4,3}^{2} v_{1,4}
  v_{4,3}^{2} +u_{4,3}^{4} +u_{4,3}^{3} v_{4,3} -u_{4,3}^{2} v_{1,4} v_{4,3}
  -2 \hspace{0.17em} u_{4,3}^{2} v_{4,3}^{2} -4 \hspace{0.17em} u_{4,3}^{3} -3
  \hspace{0.17em} u_{4,3}^{2} v_{4,3} +u_{4,3} v_{4,3}^{2} +6 \hspace{0.17em}
  u_{4,3}^{2} +3 \hspace{0.17em} u_{4,3} v_{4,3} -4 \hspace{0.17em} u_{4,3}
  -v_{4,3} +1$}}
\end{itemize}
As mentioned in the above section, the efficiency of the ``cascading
saturations'' of step [B] strongly depends on the numbering of these factors.
In the present case, a favorable permutation of the indices is $\sigma ( 1 ) =2,
\sigma ( 2 ) =1$ and $\sigma ( 3 ) =3$, other permutations might drive to
infeasible computations.

We thus set $I_{Y,2} =I_{Y} + \langle P_{Y,2} \rangle$ and saturate this ideal
by $F_{Y}$, we then compute $I_{Y,1} =I_{Y} + \langle P_{Y,1} \rangle$ and
saturate this ideal by $F_{Y} \cup \{ P_{Y,2} \}$ and finally compute $I_{Y,3}
=I_{Y} + \langle P_{Y,3} \rangle$ and saturate this ideal by $F_{Y} \cup \{
P_{Y,2} ,P_{Y,3} \}$.

\subsubsection{{\bf Step C}}

For each of these three ideals, we want to eliminate one more variable. Here are the affine polynomials we use (checking each time that the leading coefficient only vanishes at forbidden points):
\begin{enumerate}
\item In $I_{Y,1}$, we find the polynomial $v_{3,4} -v_{4,3}$ and eliminate $v_{3,4}$.
\item
In $I_{Y,2}$, the polynomial $P_{Y,2}$ displayed above is affine in $u_{4,3}$ with leading coefficient $v_{4,3}$. We eliminate $u_{4,3}$.
\item In $I_{Y,3}$, we find the following polynomial: 
$$\left(v_{1,4}^{2}v_{3,4}(v_{3,4}-1)\right) u_{4,3} -v_{1,4} v_{3,4} +1.$$ We eliminate $u_{4,3}$.
\end{enumerate}
We then define the three projections $\pi_{1}$, $\pi_{2}$ and  $\pi_{3}$ and compute (algorithm [$\alpha$]) the ideals
$I_{Y_{1}} ,I_{Y_{2}}$ and $I_{Y_{3}}$ by
substituting in $I_{Y}$ and saturating by respectively $F_{Y} ,F_{Y} \cup \{
P_{Y,1} \} ,F_{Y} \cup \{ P_{Y,1} ,P_{Y,2} \}$.

\subsubsection{{\bf Step D}}

We may now compute the prime decompositions of the three ideals. The output is as follows: the above ideal
$I_{Y_{1}}$, $I_{Y_{2}}$ are not prime and $I_{Y_{3}}$ is prime. After having
removed the embedded components, it appears that $I_{Y} +I_{Y_{1}} =I_{Y_{1}
,1} \cap I_{Y_{1} ,2}$ and $I_{Y} +I_{Y_{2}} =I_{Y_{2} ,1} \cap I_{Y_{2} ,2}$. In this decomposition,
$I_{Y_{1} ,1}$ and $I_{Y_{2} ,1}$ are prime ideals of dimension $1$, whereas $I_{Y_{1} ,2}$, $I_{Y_{2} ,2}$ and $I_{Y} +I_{Y_{3}}$ are prime ideals of dimension $2$.


\subsubsection{{\bf Specific Step $1$}}
We check that $I_{Y_{1} ,1} \cap ( I_{Y} +I_{Y_{3}} ) =I_{Y_{1} ,1}$ and
$I_{Y_{2} ,1} \cap ( I_{Y} +I_{Y_{3}} )$ so that variety
$$V ( I_{Y} ) =V ( I_{Y} +I_{Y_{1} ,2} ) \cup V ( I_{Y} +I_{Y_{2} ,2}
  ) \cup V ( I_{Y} +I_{Y_{3}} )$$ is the union of $3$ $\mathbb{Q}$-irreducible
  algebraic varieties of dimension $2$. 
  
 We check that the three ideals $I_{Y_1,2}$, $I_{Y_2,2}$ and $I_{Y_3}$ are
  principal ideals each defined by a unique polyomial $H_{i}
\in \mathbb{Q} [ Y_{i} ]$. And we additionally check that the $H_i$'s are irreducible in $\mathbb{C} [ Y_{i} ]$ (algorithm [$\theta$]). Here are the polynomials $H_i$'s:
\begin{itemize}
  \item {\small{$H_{1} =u_{4,3} v_{1,4}^{2} v_{4,3}^{2} -u_{4,3}^{2} v_{1,4}
  v_{4,3} -u_{4,3} v_{1,4}^{2} v_{4,3} -u_{4,3} v_{1,4} v_{4,3}^{2}
  +u_{4,3}^{2} v_{1,4} +3 \hspace{0.17em} u_{4,3} v_{1,4} v_{4,3} -v_{1,4}
  u_{4,3} -u_{4,3} v_{4,3} -v_{1,4} v_{4,3} +v_{4,3}$}}
  
  \item {\small{$H_{2} =v_{1,4}^{5} v_{3,4}^{2} v_{4,3}^{4} -4 \hspace{0.17em}
  v_{4,3}^{3} v_{1,4}^{5} v_{3,4}^{2} -2 \hspace{0.17em} v_{1,4}^{5} v_{3,4}
  v_{4,3}^{4} -3 \hspace{0.17em} v_{4,3}^{4} v_{1,4}^{4} v_{3,4}^{2} +6
  \hspace{0.17em} v_{4,3}^{2} v_{1,4}^{5} v_{3,4}^{2} +8 \hspace{0.17em}
  v_{1,4}^{5} v_{3,4} v_{4,3}^{3} +v_{1,4}^{5} v_{4,3}^{4} +9 \hspace{0.17em}
  v_{4,3}^{3} v_{1,4}^{4} v_{3,4}^{2} +5 \hspace{0.17em} v_{1,4}^{4} v_{3,4}
  v_{4,3}^{4} +3 \hspace{0.17em} v_{1,4}^{3} v_{3,4}^{2} v_{4,3}^{4}
  -v_{1,4}^{3} v_{3,4} v_{4,3}^{5} -4 \hspace{0.17em} v_{1,4}^{5} v_{3,4}^{2}
  v_{4,3} -12 \hspace{0.17em} v_{1,4}^{5} v_{3,4} v_{4,3}^{2} -4
  \hspace{0.17em} v_{1,4}^{5} v_{4,3}^{3} -9 \hspace{0.17em} v_{3,4}^{2}
  v_{1,4}^{4} v_{4,3}^{2} -15 \hspace{0.17em} v_{3,4} v_{4,3}^{3} v_{1,4}^{4}
  -2 \hspace{0.17em} v_{1,4}^{4} v_{4,3}^{4} -5 \hspace{0.17em} v_{1,4}^{3}
  v_{3,4}^{2} v_{4,3}^{3} -v_{1,4}^{3} v_{4,3}^{4} v_{3,4} -v_{1,4}^{2}
  v_{3,4}^{2} v_{4,3}^{4} +3 \hspace{0.17em} v_{1,4}^{2} v_{4,3}^{5} v_{3,4}
  +v_{3,4}^{2} v_{1,4}^{5} +8 \hspace{0.17em} v_{1,4}^{5} v_{3,4} v_{4,3} +6
  \hspace{0.17em} v_{1,4}^{5} v_{4,3}^{2} +3 \hspace{0.17em} v_{3,4}^{2}
  v_{1,4}^{4} v_{4,3} +15 \hspace{0.17em} v_{1,4}^{4} v_{3,4} v_{4,3}^{2} +6
  \hspace{0.17em} v_{1,4}^{4} v_{4,3}^{3} +v_{1,4}^{3} v_{3,4}^{2} v_{4,3}^{2}
  +2 \hspace{0.17em} v_{4,3}^{3} v_{1,4}^{3} v_{3,4} +v_{1,4}^{3} v_{4,3}^{4}
  -v_{1,4}^{2} v_{3,4}^{2} v_{4,3}^{3} -7 \hspace{0.17em} v_{4,3}^{4}
  v_{1,4}^{2} v_{3,4} -3 \hspace{0.17em} v_{1,4} v_{3,4} v_{4,3}^{5} -2
  \hspace{0.17em} v_{1,4}^{5} v_{3,4} -4 \hspace{0.17em} v_{1,4}^{5} v_{4,3}
  -5 \hspace{0.17em} v_{4,3} v_{1,4}^{4} v_{3,4} -6 \hspace{0.17em}
  v_{1,4}^{4} v_{4,3}^{2} +v_{1,4}^{3} v_{3,4}^{2} v_{4,3} +3 \hspace{0.17em}
  v_{3,4} v_{1,4}^{3} v_{4,3}^{2} +2 \hspace{0.17em} v_{1,4}^{2} v_{3,4}^{2}
  v_{4,3}^{2} +10 \hspace{0.17em} v_{1,4}^{2} v_{3,4} v_{4,3}^{3} +v_{1,4}
  v_{3,4}^{2} v_{4,3}^{3} +7 \hspace{0.17em} v_{1,4} v_{3,4} v_{4,3}^{4}
  +v_{3,4} v_{4,3}^{5} +v_{1,4}^{5} +2 \hspace{0.17em} v_{1,4}^{4} v_{4,3} -3
  \hspace{0.17em} v_{1,4}^{3} v_{3,4} v_{4,3} -3 \hspace{0.17em} v_{1,4}^{3}
  v_{4,3}^{2} -6 \hspace{0.17em} v_{1,4}^{2} v_{3,4} v_{4,3}^{2} -2
  \hspace{0.17em} v_{1,4}^{2} v_{4,3}^{3} -6 \hspace{0.17em} v_{1,4} v_{3,4}
  v_{4,3}^{3} -2 \hspace{0.17em} v_{3,4} v_{4,3}^{4} +2 \hspace{0.17em}
  v_{1,4}^{3} v_{4,3} +2 \hspace{0.17em} v_{1,4}^{2} v_{4,3}^{2} +v_{3,4}
  v_{4,3}^{3} +v_{1,4} v_{4,3}^{2}$}}
  
  \item {\small{$H_{3} =v_{1,4}^{5} v_{3,4}^{3} v_{4,3}^{2} -v_{1,4}^{5}
  v_{3,4}^{3} v_{4,3} +v_{1,4}^{4} v_{3,4}^{5} +v_{3,4}^{4} v_{4,3}
  v_{1,4}^{4} -2 \hspace{0.17em} v_{1,4}^{4} v_{3,4}^{3} v_{4,3}^{2} -3
  \hspace{0.17em} v_{3,4}^{4} v_{1,4}^{4} -v_{3,4}^{2} v_{1,4}^{4} v_{4,3}^{2}
  -v_{3,4}^{4} v_{4,3} v_{1,4}^{3} +v_{1,4}^{3} v_{3,4}^{3} v_{4,3}^{2} +3
  \hspace{0.17em} v_{1,4}^{4} v_{3,4}^{3} +2 \hspace{0.17em} v_{3,4}^{2}
  v_{1,4}^{4} v_{4,3} -2 \hspace{0.17em} v_{1,4}^{3} v_{3,4}^{4} +2
  \hspace{0.17em} v_{1,4}^{3} v_{3,4}^{2} v_{4,3}^{2} -v_{1,4}^{4} v_{3,4}^{2}
  +4 \hspace{0.17em} v_{1,4}^{3} v_{3,4}^{3} -2 \hspace{0.17em} v_{1,4}^{3}
  v_{3,4}^{2} v_{4,3} +v_{1,4}^{2} v_{3,4}^{3} v_{4,3} -v_{1,4}^{2}
  v_{3,4}^{2} v_{4,3}^{2} -2 \hspace{0.17em} v_{1,4}^{3} v_{3,4}^{2} +3
  \hspace{0.17em} v_{1,4}^{2} v_{3,4}^{3} +v_{1,4}^{2} v_{3,4}^{2} v_{4,3} -5
  \hspace{0.17em} v_{1,4}^{2} v_{3,4}^{2} -v_{1,4}^{2} v_{3,4} v_{4,3}
  -v_{1,4} v_{3,4}^{2} v_{4,3} +2 \hspace{0.17em} v_{1,4}^{2} v_{3,4} -2
  \hspace{0.17em} v_{1,4} v_{3,4}^{2} +v_{1,4} v_{3,4} v_{4,3} +2
  \hspace{0.17em} v_{1,4} v_{3,4} +v_{3,4} -1$.}}
\end{itemize}
Showing that, in addition, the $H_{i}$ are not singular outside the forbidden points is a 
straightforward application of the algorithms $[ \alpha ]$ and $[\gamma]$: the ideal generated 
by the partial derivatives of $H_i$, once saturated by $F_{Y_i}$, is trivial.

\subsubsection{{\bf Specific Step $2$}}

We complete the computations by a direct application of algorithm $[
\alpha ]$ in order to show that $I_{Y_{1}} +I_{Y_{2}}
+I_{Y_{3}}$ is non trivial and zero-dimensional.

\subsection{Details on algorithm [$\theta$]}

In this section we give more details on the algorithm [$\theta$] which is not available
in usual computer algebra systems.

Algorithm $[ \theta ]$ (factorization in $\mathbb{C} [ X_{1} , \ldots ,X_{n}
]$) is not present in most of computer algebra systems since they
don't know, in general, how to perform exact operations in
$\mathbb{C}$. On the other hand, most computer algebra systems implements algorithms for factorizing in
$\mathbb{Q} [ X, \ldots ,X_{n} ]$ or even $\mathbb{Q} ( \alpha ) [ X_{1} ,
\ldots ,X_{n} ]$ where $\alpha$ is any algebraic number defined by a
$\mathbb{Q}$-irreducible polynomial.

A first way to avoid this problem is to use the specific form of the $H_i$'s. Indeed, all three of them are quadratic in one variable. An elementary and careful examination, which mostly reduces to computing discriminants and checking they are not squares, leads to the proof those polynomials are irreducible over $\bC$.

More generally, a method for factorizing a polynomial $P$ over $\mathbb C$ is to find a number field
$\mathbb{Q} ( \alpha )$ such that all its factors belong to $\mathbb{Q} (
\alpha ) [ X_{1} , \ldots ,X_{n} ]$. Given a suitable $\alpha$ defined by its
minimal polynomial (with coefficients in $\mathbb{Q}$) the factorization in
$\mathbb{Q} ( \alpha ) [ X_{1} , \ldots ,X_{n} ]$ can be performed by most
computer algebra systems. Some methods for finding such a suitable $\alpha$
are reviewed in \cite{Rag97}, including the algorithm implemented in our worksheet.
\begin{comment}
 : taking integers $( a_{1} , \ldots
,a_{n-1} )$ such that $( a_{1} , \ldots ,a_{n-1} , \alpha )$ is simple point of $P
( X_{1} , \ldots ,X_{n} ) \in \mathbb{Q} [ X_{1} , \ldots ,X_{n} ]$, then $P$
has an absolute factor in $Q ( \alpha ) [ X_{1} , \ldots ,X_{n} ]$. In other terms,
minimal polynomials of suitable $\alpha$ are given by any of the
$\mathbb{Q}$-irreducible factors of $P ( a_{1} , \ldots a_{n} ,X_{n} )$.
\end{comment}

\section{Eigenvalues of peripheral representations}\label{s:A-var}

Given the triangulated structure associated with the figure eight knot
complement, the meridian and the longitude are computed from the Snappea
triangulation. The corresponding diagonal entries of the holonomy matrices are \cite{FKR14}:
\[ g_{m} = \left[ \begin{array}{ccc}
     \frac{u_{12} u_{21}}{u_{34} w_{24}} & \ast & \ast \\
     0 & 1 & \ast \\
     0 & 0 & \frac{v_{42}}{u_{34}}
   \end{array} \right] \]
\begin{center}
  $g_{l} = \left[ \begin{array}{ccc}
    \frac{u_{13} u_{31} v_{14} v_{23}}{v_{31} u_{42} v_{42} u_{24}} & \ast &
    \ast\\
    0 & 1 & \ast\\
    0 & 0 & \frac{u_{31} v_{13} u_{13} v_{24}}{u_{42} v_{32} u_{24} v_{41}}
  \end{array} \right]$
\end{center}
We therefore deduce the eigenvalues of the holonomies of the longitude and
meridian:
\[ L= \frac{u_{12} u_{21}}{u_{34} v_{24}} ,L^{\ast} = \frac{v_{42}}{u_{34}}
   ,M= \frac{u_{13} u_{31} v_{14} v_{23}}{v_{31} u_{42} v_{42} u_{24}}
   ,M^{\ast} = \frac{u_{31} v_{13} u_{13} v_{24}}{u_{42} v_{32} u_{24} v_{41}}
   . \]

\subsection{$A$-variety}

The study of these eigenvalues on each component of $\mathrm{Defor}(M,\mathcal T)$ can then be done by adding the following set $L_{\{l,m\}}$ of equations to the generators of, respectively, $I_{Y_{1}} ,I_{Y_{2}}$ and $I_{Y_{3}}$:
\begin{eqnarray*}
L_{\{ l,m \}} &= &\left\{ u_{34} v_{24} L-u_{12} u_{21}  , \quad v_{31} u_{42} v_{42} u_{24} M-u_{13} u_{31} v_{14} v_{23}  ,\right. \\
&&\left.  u_{34} L^{\ast} -v_{42} 
, \quad u_{42} v_{32}
u_{24} v_{41} M^{\ast} -u_{31} v_{13} u_{13} v_{24} \right\}.
\end{eqnarray*}

Eliminating all the variables but $\{ L,L^{\star} ,M,M^{\star} \}$ by means of
Gr{\"o}bner bases and similar factorization tricks as for the computation of
$\tmop{Defor} ( M, \mathcal{T} )$, we get that :
\begin{itemize}
  \item $( I_{Y_{2}} + \langle L_{l,m} \rangle ) \cap \mathbb{Q} [
  L,L^{\star} ,M,M^{\star} ] = \langle L^{3} -M,L^{\star 3} -M^{\star}
  \rangle$ or, equivalently, that over all the points of $D_{2}$, $L^{3} =M$
  and $L^{\star 3} =M^{\star}$
  
  \item $( I_{Y_{3}} + \langle L_{l,m} \rangle ) \cap \mathbb{Q} [
  L,L^{\star} ,M,M^{\star} ] = \langle L^{3} M-1,L^{\star 3} M^{\star} -1
  \rangle$ or, equivalently, that over all the points of $D_{3}$, $L^{3} =
  \frac{1}{M}$ and $L^{\star 3} = \frac{1}{M^{\star}}$
\end{itemize}
The eigenvalues over the component $D_{1}$ are more complicated to describe.
The Gr{\"o}bner basis of $( I_{Y_{1}} + \langle L_{l,m} \rangle ) \cap
\mathbb{Q} [ L,L^{\star} ,M,M^{\star} ]$, for the Degree Reverse
Lexicographic ordering is huge (made of 141 polynomials with up to 3462
terms). This might be due to the ordering itself or to the fact that the ideal
might have many complicated embedded components which can be viewed as an
artifact of the projection onto the coordinates $[L,L^{\star} ,M,M^{\star} ]$.

Using an elimination ordering with $M>L,L^{\star} ,M^{\star}$ one gets a
simpler description showing in particular that generically (say outside a
subvariety of dimension $1$), the points of that component of the $A$-variety
are in one-to-four correspondence with those of the hypersurface $H_{1}$ of
$\mathbb{C}^{3}$.

\subsection{Where are some already known representations ?}

\label{section:example-unipotent}Some representations are already known, more
precisely the boundary-unipotent ones, i.e. those for which the images of $l$
and $m$ are unipotent. These representations were found in {\cite{FKR14}}.
Algebraically, we add the conditions $L=L^{\ast} =M=M^{\ast} =1$. Once a point in $\mathrm{Defor}(M,\mathcal T)$ is known, it is an easy task to decide whether it belongs to $D_1$, $D_2$ or $D_3$: just plug the coordinates in the polynomials defining these three components and check which one vanishes.

Here are the representations found in \cite{FKR14} and where they live:
\begin{itemize}
  \item The monodromy of the complete hyperbolic structure on the complement of the figure-eight knot and its complex conjugate representation correspond to points of $D_1$. Indeed, if
  $\omega^{\pm} = \frac{1 \pm i \sqrt{3}}{2}$ is one root of $f_{1}$, then
  \[ u_{12} =u_{21} =u_{34} =u_{43} =v_{12} =v_{21} =v_{34} =v_{43} =
     \omega^{\pm} \]
  is a point of $\mathrm{Defor}(M,\mathcal T)$, corresponding to this monodromy or its complex conjugate. In other terms, $D_1$ is the geometric component of $\mathrm{Defor}(M,\mathcal T)$.
  
  \item Setting:
  \[ u_{1,2} =u_{3,4} =v_{3,4} =v_{4,3} = \bar{u}_{2,1} = \bar{u}_{4,3} =
     \bar{v}_{1,2} = \bar{v}_{2,1} = \omega^{\pm} , \]
     one gets two other points in $\mathrm{Defor}(M,\mathcal T)$.
  They correspond to a discrete representation of the fundamental
  group of the complement of knot in $\tmop{PU} (2,1)$ with faithful boundary
  holonomy. Moreover, its action on complex hyperbolic space has limit set the
  full boundary sphere {\cite{falbeleight}}. Both these point also belong to the geometric component $D_1$.
  
  
  
  \item Let $\gamma^{\pm} =-
  \frac{1}{2} \pm i \frac{1}{2}  \sqrt{7}$. Define
  \begin{center}
    $\begin{array}{c}
      u_{2,1} =v_{2,1} = \bar{u}_{4,3} = \bar{v}_{1,2} = \frac{5\pm i
      \sqrt{7}}{4} ,\\
      u_{1,2} =v_{2,1} = \frac{3\pm i \sqrt{7}}{8} ,v_{4,3} = \bar{v}_{3,4} =-
      \frac{1\pm i \sqrt{7}}{2}
    \end{array}$
  \end{center}
  These two points belong to $\mathrm{Defor}(M,\mathcal T)$ and correspond to representation in $\tmop{PU} (2,1)$. It turns out that they belong to $D_3$.
  \item Consider now the inverse of $\gamma^{\pm}$: $1/ \gamma^{\pm} =-
  \frac{1}{4} \pm i \frac{1}{4}  \sqrt{7}$,nd define:
  \begin{center}
    $\begin{array}{c}
      u_{1,2} = \bar{u}_{3,4} = \frac{3\pm i \sqrt{7}}{2} ,v_{4,3} =
      \bar{v}_{3,4} = \frac{5\pm i \sqrt{7}}{8} ,\\
      u_{2,1} =v_{2,1} = \bar{u}_{4,3} = \bar{v}_{1,2} = \frac{-1\pm i
      \sqrt{7}}{4}
    \end{array}$
  \end{center}  
  These two points belong to $\mathrm{Defor}(M,\mathcal T)$ and correspond to representation in $\tmop{PU} (2,1)$. It turns out that they belong to $D_2$.
  
  
\end{itemize}
All solutions were already obtained in {\cite{falbeleight}}. Moreover the solutions that belong to $D_2$ or $D_3$ correspond to spherical CR structures with
unipotent boundary holonomy of rank one {\cite{DF}}.

Moreover, in {\cite{DF}}, it was shown that there exists a non-inner
automorphism $\tau$ of the group $\Gamma_{8}$ that sends a representation
appearing in the third point above to a representation appearing in the fourth point.
It is known that, on the hyperbolic structure, this automorphism
correspond to a change of orientation, that is it acts as complex conjugation. It proves that the action of $\tau$
on the representation variety exchange the last two components and preserves
the first one. This fact is presented in a more concrete way in the following section
\ref{deformation}.

\subsubsection{Hyperbolic solutions - $A$-polynomial}

We may now look after the representations coming from $\mathrm{PGL}
(2,\mathbb{C})$, which we call hyperbolic solutions. Indeed, the hyperbolic
conditions : $u_{12} =u_{21} =u_{34} =u_{43}$, $v_{12} =v_{21} =v_{34}
=v_{43}$ on the parameters is equivalent to the fact that the holonomy
representation associated to these parameters has its image lying in
$\mathrm{Ad} ( \mathrm{PGL} (2,\mathbb{C}))$. These conditions simplify a lot
the equations and give the 1-dimensional complex variety
\[ z= \frac{1}{1-y} , \hspace{1em} x^{2} y^{2} = (1-x)  (1-y) . \]
We then obtain
\[ \begin{array}{l}
     u_{12} =u_{34} =u_{21} =u_{43} =x,\\
     u_{23} =u_{41} =u_{14} =u_{32} =1- \frac{1}{x} ,\\
     u_{13} =u_{31} =u_{24} =u_{42} = \frac{1}{1-x} ,\\
     v_{12} =v_{21} =v_{34} =v_{43} = \frac{1}{1-y} ,\\
     v_{13} =v_{24} =v_{31} =v_{42} =1- \frac{1}{y}\\
     v_{23} =v_{14} =v_{32} =v_{41} =y.
   \end{array} \]
Furthermore, we get
\[ L= \frac{1}{L^{\ast}} = \frac{xy}{x-1} ,M= \frac{1}{M^{\ast}} =
   \frac{(1-y)^{2}}{y^{4}} . \]
It appears that all hyperbolic solutions lie in the first component.

Eliminating $x$ and $y$ between $L$ and $M$ with the condition $x^{2} y^{2} =
(1-x)  (1-y)$ gives the condition (called $A$-polynomial):

{\small{$-L^{8} M+2 \hspace{0.17em} L^{7} M+3 \hspace{0.17em} L^{6} M-2
\hspace{0.17em} L^{5} M+L^{4} M^{2} -6 \hspace{0.17em} L^{4} M+L^{4} -2
\hspace{0.17em} L^{3} M+3 \hspace{0.17em} L^{2} M+2 \hspace{0.17em} LM-M=0$}}

Note that this $A$-polynomial is not the classical one for the group
$\tmop{SL} (2,\mathbb{C})$ but is the same as the one independently found by
C. Zickert {\cite{Zickert-A-pol}}.

Note that hyperbolic and CR-solutions in the first component intersect in a
1-dimensional variety, namely the real part of the hyperbolic solutions. Those
are, from a geometric point of view, degenerate as the group $\tmop{PGL}
(2,\mathbb{R})$ is not linked to a 3-dimensional geometry.
